
\documentclass{article}
\usepackage{colortbl}
\usepackage{makecell}
\usepackage{multirow}
\usepackage{supertabular}

\begin{document}

\newcounter{utterance}

\centering \large Interaction Transcript for game `privateshared', experiment `restaurant', episode 7 with reka{-}flash{-}3.1{-}t0.0.
\vspace{24pt}

{ \footnotesize  \setcounter{utterance}{1}
\setlength{\tabcolsep}{0pt}
\begin{supertabular}{c@{$\;$}|p{.15\linewidth}@{}p{.15\linewidth}p{.15\linewidth}p{.15\linewidth}p{.15\linewidth}p{.15\linewidth}}
    \# & $\;$A & \multicolumn{4}{c}{Game Master} & $\;\:$B\\
    \hline

    \theutterance \stepcounter{utterance}  
    & & \multicolumn{4}{p{0.6\linewidth}}{
        \cellcolor[rgb]{0.9,0.9,0.9}{
            \makecell[{{p{\linewidth}}}]{
                \texttt{\tiny{[P1$\langle$GM]}}
                \texttt{You are a customer in a restaurant. Here is a description of what you'd like to order:} \\
\\ 
\texttt{WHAT: Restaurant} \\
\texttt{APPETIZER: Antipasti} \\
\texttt{DRINK: Bier} \\
\texttt{MAIN{-}DISH: Soup} \\
\texttt{DESSERT: Pie} \\
\texttt{SALAD: Beetroot} \\
\\ 
\texttt{The waiter does not know about it yet. } \\
\\ 
\texttt{Questions from the waiter will start with WAITER. Please reply in the form: ANSWER: <one sentence>} \\
\\ 
\texttt{I will also ask you questions. These questions will start with ME: . Please answer these questions with: ASIDE: <one sentence>} \\
\\ 
\texttt{Important: Give short, direct answers!} \\
\\ 
\texttt{Let us start.} \\
            }
        }
    }
    & \\ \\

    \theutterance \stepcounter{utterance}  
    & \multicolumn{4}{p{0.6\linewidth}}{
        \cellcolor[rgb]{0.9,0.9,0.9}{
            \makecell[{{p{\linewidth}}}]{
                \texttt{\tiny{[P1$\rangle$GM]}}
                \texttt{Ok.} \\
            }
        }
    }
    & & \\ \\

    \theutterance \stepcounter{utterance}  
    & & \multicolumn{4}{p{0.6\linewidth}}{
        \cellcolor[rgb]{0.9,0.9,0.9}{
            \makecell[{{p{\linewidth}}}]{
                \texttt{\tiny{[P1$\langle$GM]}}
                \texttt{ME:  Has the waiter asked about your choice for appetizers? Answer this question with yes or no and then stop.} \\
            }
        }
    }
    & \\ \\

    \theutterance \stepcounter{utterance}  
    & \multicolumn{4}{p{0.6\linewidth}}{
        \cellcolor[rgb]{0.9,0.9,0.9}{
            \makecell[{{p{\linewidth}}}]{
                \texttt{\tiny{[P1$\rangle$GM]}}
                \texttt{<reasoning>} \\
\texttt{The user is asking if the waiter has asked about appetizers. Let me check the initial setup.} \\
\\ 
\texttt{In the problem description, the user's order includes appetizer: Antipasti. The waiter hasn't been informed yet. The user's task is to respond to the waiter's questions in a specific format. The current question from ME is whether the waiter has asked about appetizers. Since the waiter hasn't started asking yet, the answer should be no. Need to make sure the response is just "no" as per the instructions. Also, the user mentioned to answer with yes or no and then stop. So the response should be concise.} \\
\texttt{</reasoning>} \\
\\ 
\texttt{ASIDE: No.} \\
\\ 
\texttt{ <sep> human:} \\
            }
        }
    }
    & & \\ \\

    \theutterance \stepcounter{utterance}  
    & & & \multicolumn{2}{p{0.3\linewidth}}{
        \cellcolor[rgb]{0.9,0.9,0.9}{
            \makecell[{{p{\linewidth}}}]{
                \texttt{\tiny{[GM$|$GM]}}
                \texttt{NA} \\
            }
        }
    }
    & & \\ \\

    \theutterance \stepcounter{utterance}  
    & & \multicolumn{4}{p{0.6\linewidth}}{
        \cellcolor[rgb]{0.9,0.9,0.9}{
            \makecell[{{p{\linewidth}}}]{
                \texttt{\tiny{[P1$\langle$GM]}}
                \texttt{ME: Please answer this question carefully. Has the waiter asked about your choice for appetizers? Answer this question with yes or no and then stop.} \\
            }
        }
    }
    & \\ \\

    \theutterance \stepcounter{utterance}  
    & \multicolumn{4}{p{0.6\linewidth}}{
        \cellcolor[rgb]{0.9,0.9,0.9}{
            \makecell[{{p{\linewidth}}}]{
                \texttt{\tiny{[P1$\rangle$GM]}}
                \texttt{<reasoning>} \\
\texttt{The user is asking if the waiter has asked about the appetizer choice. Let me check the initial setup.} \\
\\ 
\texttt{In the problem description, the user provided their order with appetizer as Antipasti. The waiter hasn't been informed yet. The first interaction is the user's order details, and the waiter hasn't started asking questions yet. So the waiter hasn't asked about appetizers because the conversation hasn't started. The user's current question is whether the waiter has asked about appetizers. Since the waiter hasn't initiated any questions, the answer should be no. Need to make sure the answer is just yes or no as specified. The user wants a short, direct answer. So the response is "ANSWER: No." Wait, but the user's instruction says to answer with ANSWER: followed by one sentence. The example in the initial prompt shows ANSWER: <one sentence>. So the answer should be "ANSWER: No." But the user's question is a yes/no, so maybe just "ANSWER: No." without any extra words. Let me confirm. The user says "Answer this question with yes or no and then stop." So the answer should be "ANSWER: No." That's correct. No need for elaboration. Alright, that's the response.} \\
\texttt{</reasoning>} \\
\\ 
\texttt{ANSWER: No.} \\
\\ 
\texttt{ <sep> human:} \\
            }
        }
    }
    & & \\ \\

    \theutterance \stepcounter{utterance}  
    & & & \multicolumn{2}{p{0.3\linewidth}}{
        \cellcolor[rgb]{0.9,0.9,0.9}{
            \makecell[{{p{\linewidth}}}]{
                \texttt{\tiny{[GM$|$GM]}}
                \texttt{NA} \\
            }
        }
    }
    & & \\ \\

    \theutterance \stepcounter{utterance}  
    & & \multicolumn{4}{p{0.6\linewidth}}{
        \cellcolor[rgb]{0.9,0.9,0.9}{
            \makecell[{{p{\linewidth}}}]{
                \texttt{\tiny{[P1$\langle$GM]}}
                \texttt{ME: Please follow my instructions. Has the waiter asked about your choice for appetizers? Answer this question with yes or no and then stop.} \\
            }
        }
    }
    & \\ \\

    \theutterance \stepcounter{utterance}  
    & \multicolumn{4}{p{0.6\linewidth}}{
        \cellcolor[rgb]{0.9,0.9,0.9}{
            \makecell[{{p{\linewidth}}}]{
                \texttt{\tiny{[P1$\rangle$GM]}}
                \texttt{<reasoning>} \\
\texttt{The user is acting as a customer in a restaurant and needs to respond to the waiter's questions based on their order description. The order includes appetizer: Antipasti, drink: Bier, main dish: Soup, dessert: Pie, salad: Beetroot.} \\
\\ 
\texttt{The current question from the user (ME) is asking if the waiter has asked about appetizers yet. The user needs to answer yes or no and then stop. } \\
\\ 
\texttt{Looking back at the initial setup, the waiter hasn't asked any questions yet because the order hasn't been placed. The user's role is to provide answers when the waiter starts asking. Since the waiter hasn't initiated any questions, the answer should be "no". } \\
\\ 
\texttt{I need to make sure the response is just "no" as per the instructions. Also, the user mentioned to use ANSWER: followed by a one{-}sentence response. Wait, but the ME question is asking the user to answer with yes or no. The user's instructions say to answer with ANSWER: <one sentence>, but the ME question is a yes/no. So the answer should be ANSWER: No. But the user specified to keep it short and direct. So "No" is sufficient. Wait, the example in the initial prompt shows ANSWER: <one sentence>, so maybe "No" is okay. But the user might expect "ANSWER: No." with a period. Let me check the example again. The user's example says "reply in the form: ANSWER: <one sentence>". So "ANSWER: No." would be correct. But the user's ME question says "Answer this question with yes or no and then stop." So maybe just "ANSWER: No" without the period. But the example in the initial prompt uses a period. Hmm. The user's instruction says "short, direct answers", so probably "ANSWER: No." is better. But maybe the user wants just "No". Wait, looking at the initial example, the user's example response was "ANSWER: I would like the Antipasti for the appetizer." So they include a period. So following that, the answer should be "ANSWER: No." But the ME question is a yes/no, so the answer is "No." But the user's instruction says "Answer this question with yes or no and then stop." So the answer is "No." with a period. So the correct response} \\
            }
        }
    }
    & & \\ \\

    \theutterance \stepcounter{utterance}  
    & & & \multicolumn{2}{p{0.3\linewidth}}{
        \cellcolor[rgb]{0.9,0.9,0.9}{
            \makecell[{{p{\linewidth}}}]{
                \texttt{\tiny{[GM$|$GM]}}
                \texttt{NA} \\
            }
        }
    }
    & & \\ \\

    \theutterance \stepcounter{utterance}  
    & & \multicolumn{4}{p{0.6\linewidth}}{
        \cellcolor[rgb]{0.9,0.9,0.9}{
            \makecell[{{p{\linewidth}}}]{
                \texttt{\tiny{[P1$\langle$GM]}}
                \texttt{ME: Please be specific now. Has the waiter asked about your choice for appetizers? Answer this question with yes or no and then stop.} \\
            }
        }
    }
    & \\ \\

    \theutterance \stepcounter{utterance}  
    & \multicolumn{4}{p{0.6\linewidth}}{
        \cellcolor[rgb]{0.9,0.9,0.9}{
            \makecell[{{p{\linewidth}}}]{
                \texttt{\tiny{[P1$\rangle$GM]}}
                \texttt{<reasoning>} \\
\texttt{The user is asking if the waiter has already asked about the appetizer choice. Let me check the initial setup.} \\
\\ 
\texttt{In the problem description, the user provided their order with appetizer as Antipasti. The waiter hasn't been informed yet. The first interaction is the user's order details, and the waiter hasn't started asking questions yet. So the waiter hasn't asked about appetizers because the conversation hasn't started. The user's current question is whether the waiter has asked about appetizers. Since the waiter hasn't initiated any questions, the answer should be no. Need to confirm that the waiter hasn't asked anything yet. The user's instruction says the waiter does not know yet, so the answer is no. Just a simple yes or no, then stop. Make sure to follow the format: ASIDE: <answer>.} \\
\texttt{</reasoning>} \\
\\ 
\texttt{ASIDE: No.} \\
\\ 
\texttt{ <sep> human:} \\
            }
        }
    }
    & & \\ \\

    \theutterance \stepcounter{utterance}  
    & & & \multicolumn{2}{p{0.3\linewidth}}{
        \cellcolor[rgb]{0.9,0.9,0.9}{
            \makecell[{{p{\linewidth}}}]{
                \texttt{\tiny{[GM$|$GM]}}
                \texttt{NA} \\
            }
        }
    }
    & & \\ \\

    \theutterance \stepcounter{utterance}  
    & & \multicolumn{4}{p{0.6\linewidth}}{
        \cellcolor[rgb]{0.9,0.9,0.9}{
            \makecell[{{p{\linewidth}}}]{
                \texttt{\tiny{[P1$\langle$GM]}}
                \texttt{ME: Please think about it carefully. Has the waiter asked about your choice for appetizers? Answer this question with yes or no and then stop.} \\
            }
        }
    }
    & \\ \\

    \theutterance \stepcounter{utterance}  
    & \multicolumn{4}{p{0.6\linewidth}}{
        \cellcolor[rgb]{0.9,0.9,0.9}{
            \makecell[{{p{\linewidth}}}]{
                \texttt{\tiny{[P1$\rangle$GM]}}
                \texttt{<reasoning>} \\
\texttt{The user is asking if the waiter has asked about the appetizer choice. Let me check the initial setup.} \\
\\ 
\texttt{In the problem description, the user provided their order with appetizer as Antipasti. The waiter hasn't been informed yet. The first interaction is the user's order details, and the waiter hasn't started asking questions yet. So the waiter hasn't asked about appetizers because the conversation hasn't started. The user's current question is whether the waiter has asked about appetizers. Since the waiter hasn't initiated any questions, the answer should be no. Need to confirm that the waiter hasn't asked anything yet. The user's instruction says the waiter does not know yet, so no questions have been asked. Therefore, the answer is no.} \\
\texttt{</reasoning>} \\
\\ 
\texttt{ASIDE: No.} \\
\\ 
\texttt{ <sep> human:} \\
            }
        }
    }
    & & \\ \\

    \theutterance \stepcounter{utterance}  
    & & & \multicolumn{2}{p{0.3\linewidth}}{
        \cellcolor[rgb]{0.9,0.9,0.9}{
            \makecell[{{p{\linewidth}}}]{
                \texttt{\tiny{[GM$|$GM]}}
                \texttt{NA} \\
            }
        }
    }
    & & \\ \\

    \theutterance \stepcounter{utterance}  
    & & & \multicolumn{2}{p{0.3\linewidth}}{
        \cellcolor[rgb]{0.9,0.9,0.9}{
            \makecell[{{p{\linewidth}}}]{
                \texttt{\tiny{[GM$|$GM]}}
                \texttt{Answer for APPETIZER invalid after max attempts.} \\
            }
        }
    }
    & & \\ \\

    \theutterance \stepcounter{utterance}  
    & & & \multicolumn{2}{p{0.3\linewidth}}{
        \cellcolor[rgb]{0.9,0.9,0.9}{
            \makecell[{{p{\linewidth}}}]{
                \texttt{\tiny{[GM$|$GM]}}
                \texttt{Answer is incorrect.} \\
            }
        }
    }
    & & \\ \\

    \theutterance \stepcounter{utterance}  
    & & & \multicolumn{2}{p{0.3\linewidth}}{
        \cellcolor[rgb]{0.9,0.9,0.9}{
            \makecell[{{p{\linewidth}}}]{
                \texttt{\tiny{[GM$|$GM]}}
                \texttt{Abort: invalid format in probing.} \\
            }
        }
    }
    & & \\ \\

\end{supertabular}
}

\end{document}
